% =====================================================================
%  CARNET D'ENTRAÎNEMENT — CHIMIE — FICHE 3 — THÈME 1
%  Particules subatomiques — NIVEAU MOYEN — ÉNONCÉS
% =====================================================================
\documentclass[11pt,a4paper]{article}

\usepackage[utf8]{inputenc}
\usepackage[T1]{fontenc}
\usepackage{lmodern}
\usepackage[french]{babel}

\usepackage{amsmath,amssymb,mathtools}
\usepackage{xcolor}
\usepackage{colortbl}
\usepackage{tikz}
\usetikzlibrary{arrows.meta,positioning,calc}
\usepackage[most]{tcolorbox}
\usepackage{enumitem}
\usepackage{booktabs}
\usepackage{array}
\usepackage[a4paper,margin=2.1cm,headheight=26pt,headsep=9pt]{geometry}
\usepackage{fancyhdr}
\usepackage[hidelinks]{hyperref}

\definecolor{primary}{RGB}{124,78,140}
\definecolor{primarydark}{RGB}{74,44,92}
\definecolor{defcol}{RGB}{0,114,178}
\definecolor{thmcol}{RGB}{0,140,120}
\definecolor{attncol}{RGB}{200,40,55}

\newcommand{\nuc}[3]{\prescript{#1}{#2}{\mathrm{#3}}}
\newcommand{\pp}{p^{+}}
\newcommand{\ee}{e^{-}}
\newcommand{\nz}{n^{0}}

\pagestyle{fancy}
\fancyhf{}
\fancyhead[L]{\small\textcolor{primarydark}{\textbf{Fiche 3 — Thème 1}} : Particules subatomiques}
\fancyhead[R]{\small\textcolor{primarydark}{\textsc{Moyen — Énoncés}}}
\fancyfoot[C]{\small\thepage}
\renewcommand{\headrulewidth}{0.8pt}
\renewcommand{\headrule}{\hbox to\headwidth{\color{primary}\leaders\hrule height \headrulewidth\hfill}}

\newcounter{exo}
\newcommand{\exo}[1][]{\stepcounter{exo}\par\medskip%
  \noindent{\color{primary}\rule[-2pt]{3.5pt}{15pt}}\ \textbf{\color{primarydark}\large Exercice \theexo}%
  \ifx\relax#1\relax\relax\else\ \textmd{\normalsize\color{black!65}— \itshape #1}\fi\par\nopagebreak\smallskip}

\setlist{leftmargin=1.7em,topsep=2pt,itemsep=3pt}

\begin{document}

\begin{tikzpicture}
\node[rectangle,rounded corners=6pt,fill=primary,text=white,
      minimum width=\textwidth,minimum height=1.35cm,align=center,inner sep=6pt]
  {\Large\bfseries Thème 1 — Particules subatomiques\\[2pt]
   \normalsize\mdseries Niveau Moyen — Énoncés};
\end{tikzpicture}

\vspace{3pt}
{\small\itshape Données : $e = 1{,}6\times10^{-19}$ C ; $m(\pp)\approx m(\nz)\approx 1{,}7\times10^{-27}$ kg ;
$m(\ee)\approx 9{,}1\times10^{-31}$ kg. Travaille les puissances de $10$ à la main (multiplier les
mantisses, additionner les exposants).}

% ---------------------------------------------------------------
\exo[charge d'un grand nombre de protons]
Quelle est la charge électrique, en coulombs, d'un ensemble de $10^{12}$ protons ? Détaille le calcul en
séparant la mantisse et la puissance de $10$.

% ---------------------------------------------------------------
\exo[charge d'un ensemble d'électrons]
Un nuage contient $2{,}5\times10^{15}$ électrons. Quelle est sa charge totale, en coulombs ? Précise le
\textbf{signe} et justifie-le.

% ---------------------------------------------------------------
\exo[remonter au nombre de particules]
Un petit corps porte une charge de $+4{,}8\times10^{-18}$ C, due à un \textbf{excès de protons} (il a plus
de protons que d'électrons). Combien de protons en excès porte-t-il ?

% ---------------------------------------------------------------
\exo[rapport des masses]
\begin{enumerate}[label=\alph*)]
  \item Calcule le rapport $\dfrac{m(\ee)}{m(\pp)}$. Combien de fois, environ, l'électron est-il plus léger
        que le proton ?
  \item Un proton est-il plus lourd, aussi lourd, ou plus léger qu'un neutron ?
\end{enumerate}

% ---------------------------------------------------------------
\exo[où se cache la masse d'un atome ?]
On considère un atome dont le noyau contient $6$ protons et $6$ neutrons ; l'atome neutre possède donc
$6$ électrons.
\begin{enumerate}[label=\alph*)]
  \item Estime la masse du \textbf{noyau} (somme des masses des nucléons), en kg.
  \item Estime la masse totale des $6$ \textbf{électrons}, en kg.
  \item Compare : quelle fraction de la masse de l'atome les électrons représentent-ils, en ordre de
        grandeur ? Conclus.
\end{enumerate}

\end{document}
